\section{Benutzererlebnis}

In diesem Entwicklungsbereich wurden die äußere Gestaltung und die wahrnehmbare Interaktion von Pablo ausgearbeitet. 
Ziel war es, die technischen Funktionen in eine konsistente und verständliche Produktgestalt zu überführen.

\subsection{Ergebnis}

Das Gehäuse von Pablo wurde als 3D-gedruckter Baumstamm ausgeführt und integriert die relevanten Hardwarekomponenten des Systems. 
Im Kopfbereich befindet sich ein Display, auf dem Pablo als animiertes Pixel-Art-Faultier dargestellt wird. 
Die Darstellung umfasst unterschiedliche visuelle Zustände in Abhängigkeit vom Systemstatus. 
Zusätzlich wurde eine mechanische Kopfbewegung umgesetzt, bei der sich Pablo nach der Aktivierung in die Richtung der sprechenden Person ausrichtet.

\subsection{Einordnung}

Die Gestaltung des Benutzererlebnisses dient dazu, Pablo nicht nur als technische Funktionseinheit, sondern als physisch wahrnehmbares Interaktionssystem auszubilden. 
Die visuelle Darstellung als Faultier und die Einbettung in ein baumstammförmiges Gehäuse schaffen eine wiedererkennbare und weniger technisch wirkende Produktgestalt. 
Die gerichtete Kopfbewegung ergänzt dies um eine sichtbare Reaktion auf die Ansprache und macht den Interaktionsbeginn für Nutzende unmittelbar nachvollziehbar. 
Insgesamt unterstützen diese Elemente eine klare, konsistente und niedrigschwellige Mensch-Maschine-Interaktion.